% ch4_d 第4章 §4.3尾–§4.4 (书页 177–186 / p192–p201)
取 $\omega\to 0$，我们得到下面这个非常有用的直流电导率公式：
\begin{equation*}
\sigma^{ij}(0) = \int \frac{\mathrm{d}\nu}{\pi}\, \frac{\mathrm{d}^d\pmb{k}}{(2\pi)^d}\,
\frac{\Gamma^2 v^i(\pmb{k}) v^j(\pmb{k}) \left(-\dfrac{\partial n_F(\nu)}{\partial \nu}\right)}
{\left((\nu - \xi_{\pmb{k}})^2 + \Gamma^2\right)^2}
\approx \int \frac{\mathrm{d}^d\pmb{k}}{(2\pi)^d}\, \tau\, v^i(\pmb{k}) v^j(\pmb{k})\, \delta(\xi_{\pmb{k}})
\end{equation*}
其中我们假定了 $T$ 与 $\Gamma$ 都很小，并利用了 $-\partial n_F(\nu)/\partial\nu \approx \delta(\nu)$
以及 $\Gamma^2/(\xi_{\pmb{k}}^2 + \Gamma^2)^2 \approx \frac{\pi}{2\Gamma}\,\delta(\xi_{\pmb{k}})$。

\subsection{其他两体关联函数}

当电子具有自旋时，我们还可以计算自旋关联函数和电子对关联函数。自旋密度算符由
$s^i = c_\alpha^\dagger \sigma^i_{\alpha\beta} c_\beta/2$ 给出，其中 $c_\beta$（$\beta = 1, 2$）
为自旋向上与自旋向下的电子算符，$\sigma^i$ 为泡利矩阵。自旋--自旋关联等于密度关联：
\begin{equation*}
\big\langle s^3(\pmb{x},t) s^3(\pmb{0}) \big\rangle
= \frac{1}{4} \Big[ \big\langle (c_1^\dagger c_1)(c_1^\dagger c_1) \big\rangle
+ \big\langle (c_2^\dagger c_2)(c_2^\dagger c_2) \big\rangle \Big]
= \frac{1}{2\mathrm{i}} \big( \Pi^{00}(\pmb{x},t) - \rho_0^2 \big)
\end{equation*}
我们看到自旋--自旋关联也具有代数长程序。对于嵌套费米面，
\begin{equation*}
\big\langle s^i(\pmb{x},0) s^j(\pmb{0}) \big\rangle \propto \delta_{ij}\, \frac{1 + \sin(Q|\pmb{x}|)}{|\pmb{x}|^2}
\end{equation*}
若 $\pmb{x} \| \pmb{Q}$。

电子对算符由 $b(\pmb{x},t) = c_1(\pmb{x},t)c_2(\pmb{x},t)$ 给出，它与超导电性有关。其关联等于无自旋电子的格林函数的平方：
\begin{equation*}
\big\langle b(\pmb{x},t) b^\dagger(\pmb{0}) \big\rangle = G^2(\pmb{x},t)
\propto \Big|_{t=0} \cos^2 \left( k_F |\pmb{x}| - \frac{\pi(d+1)}{4} \right) \frac{1}{|\pmb{x}|^{(d+1)}}
\end{equation*}
在电子对关联中我们再一次得到代数长程序。

总结起来，我们发现自由费米子系统中的许多关联都具有幂律衰减。这表明自由费米子基态是一种临界态。有时，任意弱的相互作用就能把代数长程序变成真正的长程序，并诱导自发对称性破缺。例如，无限弱的吸引相互作用就会在电子对关联中产生长程序，并生成超导态。

%FIG{4.16}: {$\dfrac{1}{(z-\alpha-\mathrm{i}0^{+})^{2}+1-(\alpha+\mathrm{i}0^{+})^{2}}$ 的两个极点：(a) $|\alpha|>1$ 情形；(b) $|\alpha|<1$ 情形。若不存在 $\mathrm{i}0^{+}$ 项，$|\alpha|<1$ 的两个极点将位于单位圆 $|z|=1$ 上。}

\subsection{注记：一些计算细节}

\subsubsection{二维中 $\pi^{\perp}$ 的计算}

我们如下计算 $\pi^{22}$：
\begin{align*}
\pi^{22}(k\hat{\pmb{x}},\omega)
&= \frac{1}{4m^{2}}\, \frac{1}{4\pi^{2}} \int \mathrm{d}\Big(\frac{q^{2}}{2m}\Big) \mathrm{d}\theta\;
\delta(\epsilon(\pmb{q}) - \mu)\, \frac{\pmb{k}\cdot\pmb{q}}{\omega - \frac{\pmb{k}\cdot\pmb{q}}{m} + \mathrm{i}0^{+}}\; (2q_{2} + k_{2})^{2}\\
&= \frac{1}{m^{2}}\, \frac{1}{4\pi^{2}} \int \mathrm{d}\theta\;
\frac{k k_{F} \cos(\theta)\, k_{F}^{2} \sin^{2}(\theta)}{\omega - \frac{k k_{F}\cos(\theta)}{m} + \mathrm{i}0^{+}}
= \frac{k_{F}^{2}}{4\pi^{2} m} \int \mathrm{d}\theta\;
\frac{\cos(\theta) \sin^{2}(\theta)}{\frac{\omega}{v_{F} k} - \cos(\theta) + \mathrm{i}0^{+}}
\end{align*}
以及（见图 4.16）
\begin{align*}
\int \mathrm{d}\theta\; \frac{\cos(\theta) \sin^{2}(\theta)}{\alpha + \mathrm{i}0^{+} - \cos(\theta)}
&= \frac{1}{4} \oint_{|z|=1} \mathrm{d}z\; \frac{1}{\mathrm{i} z^{3}}\, \frac{(z^{2}-1)^{2}(z^{2}+1)}{(z-\alpha)^{2} + 1 - \alpha^{2}}\\
&= \left\{
\begin{array}{ll}
-\pi + 2\pi\alpha^{2} - 2\pi \mathrm{i} \alpha \sqrt{1-\alpha^{2}} & \text{对于 } |\alpha| < 1\\
-\pi + 2\pi\alpha^{2} - 2\pi |\alpha| \sqrt{\alpha^{2}-1} & \text{对于 } |\alpha| > 1
\end{array} \right.
\end{align*}

\subsubsection{三维中 $\pi^{\perp}$ 的计算}

为计算 $\pi^{\perp}$，我们可取 $\pmb{k} = (0, 0, k)$ 并利用
\begin{align*}
\pi^{\perp}(\omega, k) &= \pi^{11}(k\hat{\pmb{z}}, \omega)
= \frac{1}{4m^{2}}\, \frac{1}{8\pi^{3}} \int q\, \mathrm{d}\Big(\frac{q^{2}}{2m}\Big) \mathrm{d}\phi\, \mathrm{d}\theta\;
\frac{\cos(\theta)\, \delta(\epsilon(\pmb{q}))\, \pmb{k}\cdot\pmb{q}}{\omega - \frac{\pmb{k}\cdot\pmb{q}}{m} + \mathrm{i}0^{+}}\; (2q_{1} + k_{1})^{2}\\
&= \frac{1}{m^{2}}\, \frac{1}{8\pi^{3}} \int \mathrm{d}\phi\, \mathrm{d}\theta\;
\frac{k \cos(\theta)\, k_{F}^{4} \sin^{3}(\theta) \cos^{2}(\phi)}{\omega - \frac{k k_{F}\cos(\theta)}{m} + \mathrm{i}0^{+}}
= \frac{k_{F}^{3}}{8\pi^{2} m} \int_{-1}^{1} \mathrm{d}t\; \frac{t(1-t^{2})}{\frac{\omega}{v_{F} k} - t + \mathrm{i}0^{+}}\\
&= \frac{k_{F}^{3}}{8\pi^{2} m} \left( -\frac{4}{3} + 2\frac{\omega^{2}}{v_{F}^{2} k^{2}}
- \frac{\omega}{v_{F} k} \Big( 1 - \frac{\omega^{2}}{v_{F}^{2} k^{2}} \Big)
\Big( \ln \frac{|\omega - v_{F} k|}{|\omega + v_{F} k|} + \mathrm{i}\pi \Theta(v_{F} k - |\omega|) \Big) \right)
\end{align*}
我们还注意到 $\rho = k_{F}^{3}/6\pi^{2}$。

\section{绝缘体的线性响应与量子化霍尔电导}

\begin{keypoints}
\item 绝缘体的霍尔电导是一个拓扑量，并且是量子化的。
\end{keypoints}

在上一节中，我们讨论了具有费米面的自由费米子系统的线性响应。这些响应是由跨费米面的粒子--空穴激发引起的。在本节中，我们想讨论绝缘体的线性响应。这里，``绝缘体''指的是任何具有有限能隙的态。（这样的态也称为刚性态（rigid state）。）由于没有低能激发，响应的虚部在小 $\omega$ 处为零（即不存在低能耗散）。

考虑一个与电磁场耦合的相互作用费米子系统：$\mathcal{L}(c, c^{\dagger}, A_{\mu}) + \mathcal{L}_{\mathrm{gauge}}(A_{\mu})$。我们将要得到的结果非常普遍，甚至不需要知道 $\mathcal{L}(c, c^{\dagger}, A_{\mu})$ 的具体形式。我们只假设费米子的基态具有有限能隙，并且 $\mathcal{L}(c, c^{\dagger}, A_{\mu})$ 在规范变换
\begin{equation*}
c(x^{\mu}) \to \mathrm{e}^{\mathrm{i}\phi(x^{\mu})} c(x^{\mu}), \qquad A_{\mu} \to A_{\mu} + \partial_{\mu} \phi(x^{\mu}).
\end{equation*}
下不变。把费米子积掉之后，我们得到规范场的一个有效作用量：
\begin{align*}
&\mathcal{L}_{\mathrm{eff}}(A_{\mu}) = \mathcal{L}_{\mathrm{gauge}}(A_{\mu}) + \delta\mathcal{L}_{\mathrm{eff}}(A_{\mu})\\
&\delta\mathcal{L}_{\mathrm{eff}}(A_{\mu}) = -\frac{1}{2} \int \mathrm{d}x\, \mathrm{d}x'\; A_{\mu}(x) P^{\mu\nu}(x - x') A_{\nu}(x')
\end{align*}
其中 $\delta\mathcal{L}_{\mathrm{eff}}(A_{\mu})$ 是费米子的贡献，$P^{\mu\nu}(x-x')$ 是流--流关联函数。

对于绝缘体（或刚性态），$\delta\mathcal{L}_{\mathrm{eff}}$ 是局域的，而且 $P^{\mu\nu}_{k_{\mu}}$ 在频率--动量空间中是 $k_{\mu}$ 的多项式。这与上一节为金属算出的 $P^{\mu\nu}_{k_{\mu}}$ 形成鲜明对比。金属在小 $(\pmb{k}, \omega)$ 极限下的奇异性是由跨费米面的无能隙粒子--空穴激发引起的。

这里 $\delta\mathcal{L}_{\mathrm{eff}}(A_{\mu})$ 也应当是规范不变的。因此，$\delta\mathcal{L}_{\mathrm{eff}}(A_{\mu})$ 应当只通过场强 $\pmb{E}$ 和 $\pmb{B}$ 依赖于 $A_{\mu}$。于是 $\delta\mathcal{L}_{\mathrm{eff}}$ 具有如下形式：
\begin{equation*}
\delta\mathcal{L}_{\mathrm{eff}} = -\frac{1}{2} \big( B_{i} \chi^{ij} B_{j} + E_{i} p^{ij} E_{j} \big) + \dots
\end{equation*}
其中``$\dots$''代表更高阶导数项。张量 $\chi^{ij}$ 是磁化率，$p^{ij}$ 代表介电常数的一个修正。一般而言，以上就是绝缘体线性响应仅有的几种形式。

然而，在 $2+1$ 维中，$\delta\mathcal{L}_{\mathrm{eff}}$ 还可以包含一个新项——Chern--Simons 项——如下：
\begin{equation*}
\delta\mathcal{L}_{\mathrm{eff}} = \frac{K}{4\pi} A_{\mu} \partial_{\nu} A_{\lambda} \epsilon^{\mu\nu\lambda} - \frac{1}{2} \big( B_{i} \chi^{ij} B_{j} + E_{i} p^{ij} E_{j} \big) + \dots
\end{equation*}
其中 $\mu, \nu, \lambda = 0, 1, 2$，而 $\epsilon^{\mu\nu\lambda}$ 是全反对称张量。尽管 Chern--Simons 项无法用场强写出，它却仍然是规范不变的。在规范变换
\begin{equation*}
A_{\mu} \to A_{\mu} + \partial_{\mu} f
\end{equation*}
下，作用量发生如下变化：
\begin{equation*}
S = \int_{V} \mathrm{d}x\; \mathcal{L}_{\mathrm{eff}} \to S + \oint_{S} \mathrm{d}S_{\mu}\, f\, \frac{K}{4\pi} \partial_{\nu} A_{\lambda} \epsilon^{\mu\nu\lambda}
\end{equation*}
其中 $V$ 是时空体积，$S$ 是 $V$ 的表面。我们看到，如果我们的系统生活在无边界的闭合时空上，那么 Chern--Simons 项就是规范不变的。

Chern--Simons 项的线性响应为
\begin{equation*}
J^{i} = -\frac{\delta S}{\delta A_{i}} = \frac{K}{2\pi} \partial_{0} A_{i} \epsilon^{ij} = \frac{K}{2\pi} E_{i} \epsilon^{ij}
\end{equation*}
我们看到 Chern--Simons 项给出霍尔电导 $\sigma_{xy} = K/2\pi$（如果把 $e$ 和 $\hbar$ 放回去，则 $\sigma_{xy} = Ke^{2}/h$）。

下面我们将证明霍尔电导 $\sigma_{xy}$ 或 $K$ 是量子化的（Thouless et al., 1982; Avron et al., 1983）。首先，让我们考虑一个周期系统和一个特殊的规范场位形：
\begin{equation*}
A_{1} = \frac{\theta_{1}(t)}{L_{1}}, \qquad A_{2} = \frac{\theta_{2}(t)}{L_{2}}, \qquad A_{0} = 0 \tag{4.4.1}
\end{equation*}
其中 $L_{1,2}$ 是系统在 $x$ 与 $y$ 方向上的尺寸。费米子的多体基态 $|\pmb{\theta}\rangle$ 由 $\pmb{\theta} = (\theta_{1}, \theta_{2})$ 参量化。如果我们把费米子积掉，就会得到一个有效拉格朗日量 $\delta L_{\mathrm{eff}} = \int \mathrm{d}^{2}\pmb{x}\; \delta\mathcal{L}_{\mathrm{eff}}$：
\begin{equation*}
\delta L_{\mathrm{eff}} = \frac{K}{4\pi} \theta_{i} \dot{\theta}_{j} \epsilon^{ij} + \dot{\theta}_{i} p^{ij} \dot{\theta}_{j}
\end{equation*}
其中第一项来自 Chern--Simons 项，第二项来自 $\pmb{E}^{2}$ 项。如果我们把 $\pmb{\theta}$ 看作一个粒子的坐标，那么 $\delta S$ 描写的就是一个处于均匀``磁场'' $b$ 中的二维带电粒子：
\begin{align*}
\delta L_{\mathrm{eff}} &= a_{i}(\pmb{\theta}) \dot{\theta}_{i} + \dot{\theta}_{i} p^{ij} \dot{\theta}_{j}\\
a_{1} &= -\frac{K}{4\pi} \theta_{2}, \qquad a_{2} = \frac{K}{4\pi} \theta_{1}, \qquad b = \partial_{\theta_{1}} a_{2} - \partial_{\theta_{2}} a_{1} = \frac{K}{2\pi}
\end{align*}

在 4.4.1 节中，我们将证明费米子态 $|\theta_{1}, \theta_{2}\rangle$、$|\theta_{1} + 2\pi, \theta_{2}\rangle$ 与 $|\theta_{1}, \theta_{2} + 2\pi\rangle$ 由规范变换联系在一起。物理上，这意味着 $|\theta_{1}, \theta_{2}\rangle$、$|\theta_{1} + 2\pi, \theta_{2}\rangle$ 与 $|\theta_{1}, \theta_{2} + 2\pi\rangle$ 其实是同一个物理态（见 6.1 节的讨论）。因此，$\delta L_{\mathrm{eff}}$ 实际上描写的是一个环面上的粒子，环面由 $0 < \theta_{1} < 2\pi$ 与 $0 < \theta_{2} < 2\pi$ 参量化。

让我们让粒子沿回路 $C$：$(0,0) \to (2\pi,0) \to (2\pi,2\pi) \to (0,2\pi) \to (0,0)$ 移动。这样一圈移动所积累的相位由 $C$ 所包围的``磁''通量给出：
\begin{equation*}
\oint_{C} \mathrm{d}\pmb{\theta} \cdot \pmb{a} = \int_{D} \mathrm{d}^{2}\pmb{\theta}\; b = 2\pi K
\end{equation*}
其中 $D$ 是回路 $C$ 所围的面积。在任何闭合曲面（例如球面或环面）上，任何回路都会围住该回路两侧的两个曲面（见图 6.5）。这里的 $D$ 只是 $C$ 所围两个曲面中的一个。$C$ 所围的另一个曲面 $D'$ 面积为零（因为 $D$ 覆盖了环面的全部面积）。所以相位 $\oint_{C} \mathrm{d}\pmb{\theta} \cdot \pmb{a}$ 也可以写成 $\int_{D'} \mathrm{d}^{2}\pmb{\theta}\; b = 0$。为保持一致，$\int_{D} \mathrm{d}^{2}\pmb{\theta}\; b$ 与 $\int_{D'} \mathrm{d}^{2}\pmb{\theta}\; b$ 应当代表同一个相位。于是
\begin{equation*}
\int_{D} \mathrm{d}^{2}\pmb{\theta}\; b = 2\pi \times \text{整数}
\end{equation*}

从数学上可以证明，穿过任何闭合曲面的总磁通量必须量子化为整数。这正是磁单极子的磁荷量子化的原因。对我们的情形，这意味着 $K$ 量子化为整数（另见 4.4.1 节）。

我们想强调，$K$ 的量子化是非常普遍的。它适用于费米子和/或玻色子的任何相互作用系统。唯一的假设是基态 $|\pmb{\theta}\rangle$ 不简并。于是，\emph{如果环面上的多体基态不简并且具有有限能隙，那么系统的霍尔电导就量子化为整数 $\times\ e^{2}/h$}。我们想指出，如果多体基态具有简并（见 8.2.1 节），那么 $K$ 和霍尔电导就可以是一个\emph{有理}数（Niu et al., 1985）。

霍尔电导的量子化有一些有趣的推论。让我们考虑均匀磁场中的一个无相互作用电子系统。如果 $n$ 个朗道能级被填满，系统就会有有限能隙，且霍尔电导为 $ne^{2}/h$（即 $K = n$）。现在让我们加入相互作用、周期势，甚至无规势。只要能隙在此过程中始终不闭合，霍尔电导就不能改变！由于霍尔电导对任何微扰都是稳健的，我们称它为一个\emph{拓扑}量子数。改变霍尔电导的唯一办法是使能隙闭合，这将诱导一个量子相变。

为了显式计算 $K$，我们注意到 $\int \mathrm{d}t\; \delta L_{\mathrm{eff}}$ 正是绝热演化 $|\pmb{\theta}(t)\rangle$ 的作用量。于是我们有 $\int \mathrm{d}t\; \delta L_{\mathrm{eff}} = \int \mathrm{d}t\; \mathrm{i}\langle \pmb{\theta}(t) | \frac{\mathrm{d}}{\mathrm{d}t} | \pmb{\theta}(t) \rangle$，或者（见式 (2.3.4)）
\begin{equation*}
\frac{K}{4\pi} \theta_{i} \dot{\theta}_{j} \epsilon^{ij} = \mathrm{i} \Big\langle \pmb{\theta}(t) \Big| \frac{\mathrm{d}}{\mathrm{d}t} \Big| \pmb{\theta}(t) \Big\rangle \tag{4.4.2}
\end{equation*}

这里我们想说明，作为费米子的基态，$|\pmb{\theta}(t)\rangle$ 的相位是不固定的。如果我们重新定义 $|\pmb{\theta}\rangle$ 的相位：$|\pmb{\theta}\rangle \to \mathrm{e}^{\mathrm{i}\varphi(\pmb{\theta})} |\pmb{\theta}\rangle$，那么贝里相项将改变一个全导数项：
\begin{equation*}
\Big\langle \pmb{\theta}(t) \Big| \frac{\mathrm{d}}{\mathrm{d}t} \Big| \pmb{\theta}(t) \Big\rangle \to \Big\langle \pmb{\theta}(t) \Big| \frac{\mathrm{d}}{\mathrm{d}t} \Big| \pmb{\theta}(t) \Big\rangle + \mathrm{i}\, \frac{\mathrm{d}\varphi}{\mathrm{d}t}
\end{equation*}
因此，Chern--Simons 项与贝里相之间的关系式 (4.4.2) 不可能是正确的，因为开放路径的贝里相甚至没有良定义（见 2.3 节）。正确的关系应当是式 (4.4.2) 沿一条闭合回路的积分：
\begin{equation*}
\oint \mathrm{d}t\; \frac{K}{4\pi} \theta_{i} \dot{\theta}_{j} \epsilon^{ij} = \mathrm{i} \oint \mathrm{d}t\; \Big\langle \pmb{\theta}(t) \Big| \frac{\mathrm{d}}{\mathrm{d}t} \Big| \pmb{\theta}(t) \Big\rangle
\end{equation*}
特别地，沿回路 $C$，贝里相由下式给出：
\begin{equation*}
2\pi K = \frac{\sigma_{xy} \hbar}{e^{2}} = \oint_{C} \mathrm{d}\pmb{\theta} \cdot \langle \pmb{\theta} | \mathrm{i} \partial_{\pmb{\theta}} | \pmb{\theta} \rangle \tag{4.4.3}
\end{equation*}

让我们用贝里相 (4.4.3) 来计算一个能带绝缘体的量子化霍尔电导，对此系统有
\begin{equation*}
H = \sum_{ij} c_{i}^{\dagger} t_{ij} c_{j} = \sum_{\pmb{k}} c_{\pmb{k}}^{\dagger} M(\pmb{k}) c_{\pmb{k}} \tag{4.4.4}
\end{equation*}
其中 $c_{i}$ 有 $n$ 个分量，$t_{ij}$ 是只依赖 $i - j$ 的 $n \times n$ 矩阵，$M(\pmb{k})$ 是 $t_{ij}$ 的傅里叶变换。与具有式 (4.4.1) 所给形式的均匀规范势 $\pmb{A}$ 耦合后，我们得到
\begin{equation*}
H = \sum_{ij} c_{i}^{\dagger} t_{ij}\, \mathrm{e}^{\mathrm{i}\pmb{A}\cdot(\pmb{i}-\pmb{j})} c_{j} = \sum_{\pmb{k}} c_{\pmb{k}}^{\dagger} M^{\theta}(\pmb{k}) c_{\pmb{k}}
\end{equation*}
然后设 $\psi^{\theta}_{a,\pmb{k}}$ 为 $M^{\theta}(\pmb{k})$ 的本征矢。这里 $\psi^{\theta}_{a,\pmb{k}}$ 由晶体动量 $\pmb{k}$ 标记，而 $a = 1, \dots, n$，其中 $a$ 标记第 $a$ 个本征矢。注意 $\psi^{\theta}_{a,\pmb{k}}$ 是一个 $n$ 维复矢量。

让我们假设基态 $|\pmb{\theta}\rangle$ 由填充 $a = 1$ 能带得到。我们有
\begin{equation*}
| \pmb{\theta} \rangle = \otimes_{\pmb{k}} \big| \psi^{\theta}_{1,\pmb{k}} \big\rangle
\end{equation*}
利用贝里相的可加关系，即
\begin{equation*}
\oint_{C} \mathrm{d}\pmb{\theta} \cdot \langle \pmb{\theta}, 0 | \mathrm{i} \partial_{\pmb{\theta}} | \pmb{\theta}, 0 \rangle
= \oint_{C} \mathrm{d}\pmb{\theta} \cdot \langle \pmb{\theta}, 1 | \mathrm{i} \partial_{\pmb{\theta}} | \pmb{\theta}, 1 \rangle
+ \oint_{C} \mathrm{d}\pmb{\theta} \cdot \langle \pmb{\theta}, 2 | \mathrm{i} \partial_{\pmb{\theta}} | \pmb{\theta}, 2 \rangle
\end{equation*}
若 $|\pmb{\theta}, 0\rangle = |\pmb{\theta}, 1\rangle \otimes |\pmb{\theta}, 2\rangle$，我们发现
\begin{equation*}
\oint_{C} \mathrm{d}\pmb{\theta} \cdot \langle \pmb{\theta} | \mathrm{i} \partial_{\pmb{\theta}} | \pmb{\theta} \rangle
= \sum_{\pmb{k}} \oint_{C} \mathrm{d}\pmb{\theta} \cdot (\psi^{\theta}_{1,\pmb{k}})^{\dagger} \mathrm{i} \partial_{\pmb{\theta}} \psi^{\theta}_{1,\pmb{k}} \tag{4.4.5}
\end{equation*}

%FIG{4.17}: {固定的 $\pmb{k}$ 下由 $\pmb{k}' = \pmb{k} + \pmb{\theta} L^{-1}$ 扫出的小回路。所有不同 $\pmb{k}$ 的小回路覆盖整个布里渊区。沿所有小回路的回路积分对应于沿回路 $C_{\pmb{k}}$ 的回路积分，因为来自内部连线的贡献相互抵消。}

我们如何计算 $\psi^{\theta}_{1,\pmb{k}}$ 的贝里相？这里我们注意到 $M^{\theta}(\pmb{k}) = M(\pmb{k} + \pmb{\theta} L^{-1})$，其中为简单起见我们已在式 (4.4.1) 中取了 $L_{1} = L_{2} = L$。于是，$\psi^{\theta}_{1,\pmb{k}}$ 是 $M(\pmb{k} + \pmb{\theta} L^{-1})$ 的本征矢。设 $\psi_{a,\pmb{k}'}$ 为 $M(\pmb{k}')$ 的本征矢；那么当 $\pmb{k}' = \pmb{k} + \pmb{\theta} L^{-1}$ 时，$\psi^{\theta}_{1,\pmb{k}} = \psi_{1,\pmb{k}'}$。当 $\pmb{\theta}$ 沿回路 $C$ 移动时，对固定的 $\pmb{k}$，$\pmb{k}'$ 沿一个小回路移动（见图 4.17）。由于分立的 $\pmb{k}$ 具有 $(n_{1} \frac{2\pi}{L}, n_{2} \frac{2\pi}{L})$ 的形式，我们看到不同 $\pmb{k}$ 的所有小回路互不重叠地覆盖了整个布里渊区。因此，式 (4.4.5) 中各 $\psi^{\theta}_{1,\pmb{k}}$ 的贝里相之和等于 $\psi_{1,\pmb{k}'}$ 沿一条包围整个布里渊区的大回路 $C_{\pmb{k}}$ 的贝里相：
\begin{align*}
2\pi K &= \oint_{C_{\pmb{k}}} \mathrm{d}\pmb{k} \cdot \psi_{1,\pmb{k}}^{\dagger}\, \mathrm{i} \partial_{\pmb{k}} \psi_{1,\pmb{k}}\\
&= \int \mathrm{d}^{2}\pmb{k}\; \mathrm{i} \big[ (\partial_{k_{x}} \psi_{1,\pmb{k}}^{\dagger}) (\partial_{k_{y}} \psi_{1,\pmb{k}}) - (\partial_{k_{y}} \psi_{1,\pmb{k}}^{\dagger}) (\partial_{k_{x}} \psi_{1,\pmb{k}}) \big] \tag{4.4.6}
\end{align*}
如果填充了几个能带，那么我们需要对每个能带的贡献求和。

作为一个具体例子，考虑正方格子上如下的依赖自旋的跃迁哈密顿量：
\begin{equation*}
H = \sum_{i} c_{i}^{\dagger} t_{ij} c_{j}
\end{equation*}
其中 $2 \times 2$ 的跃迁矩阵为
\begin{equation*}
t_{\pmb{i}, \pmb{i}+\pmb{x}} = t \sigma^{x}, \qquad t_{\pmb{i}, \pmb{i}+\pmb{y}} = t \sigma^{y}, \qquad t_{\pmb{i}, \pmb{i}+\pmb{x}+\pmb{y}} = t' \sigma^{z},
\end{equation*}

%FIG{4.18}: {箭头表示 $\pmb{B}$ 的方向。缠绕数仅收到来自四个点 $\pmb{k} = (\pm\pi/2, \pm\pi/2)$ 附近的贡献。在每个点附近，$\pmb{B}$ 覆盖半球，贡献缠绕数 $1/2$。总缠绕数为 2。}

在动量空间中，哈密顿量变为
\begin{equation*}
H_{\pmb{k}} = 2 t \cos(k_{x}) \sigma^{x} + 2 t \cos(k_{x}) \sigma^{y} + 2 t' \cos(k_{x} + k_{y}) \sigma^{z} = -\pmb{B}(\pmb{k}) \cdot \sigma
\end{equation*}
（译注：按物理内容与图 4.18，$\sigma^{y}$ 项的系数应为 $\cos(k_{y})$；原书该项也印作 $\cos(k_{x})$。）

我们看到系统有两个能带。这里的 $H_{\pmb{k}}$ 就像磁场 $\pmb{B}(\pmb{k})$ 中一个自旋 $1/2$ 的自旋的哈密顿量。较低的能带对应于沿 $\pmb{B}$ 方向的自旋。当沿 $\pmb{k}$ 空间中的一条回路移动时，沿 $\pmb{B}(\pmb{k})$ 方向的自旋在自旋空间的单位球上扫出一条闭合回路，并产生一个贝里相。如果填充较低的能带，那么只要选取包围整个布里渊区的回路，霍尔电导就可以由上述自旋 $1/2$ 的贝里相算出。从小 $t'$ 极限下 $\pmb{B}(\pmb{k})$ 的表达式我们看到，除四个点 $\pmb{k} = (\pm\pi/2, \pm\pi/2)$ 附近外，$\pmb{B}$ 都在 $x$--$y$ 平面内。考察 $\pmb{k} = (\pm\pi/2, \pm\pi/2)$ 附近的自旋（见图 4.18），我们发现，当 $\pmb{k}$ 扫过布里渊区时，$\pmb{B}/|\pmb{B}|$ 扫过单位球两次。或者更精确地说，$\pmb{B}(\pmb{k})/|\pmb{B}(\pmb{k})|$ 以缠绕数 2 把布里渊区 $S^{1} \times S^{1}$ 映射到单位球 $S^{2}$ 上。总贝里相因此为 $2 \times 2\pi$。我们求得 $K = 2$，半填充的跃迁系统的霍尔电导为 $\sigma_{xy} = 2e^{2}/h$。

\subsection*{问题 4.4.1}

对于哈密顿量 (4.4.4)，如果 $a = 1$ 能带被部分填充，试证明霍尔电导由对被填充能级的积分给出：
\begin{equation*}
\sigma_{xy} = \frac{e^{2}}{h}\, \frac{1}{2\pi} \int_{\text{filled}} \mathrm{d}^{2}\pmb{k}\; \mathrm{i} \big[ (\partial_{k_{x}} \psi_{1,\pmb{k}}^{\dagger}) (\partial_{k_{y}} \psi_{1,\pmb{k}}) - (\partial_{k_{y}} \psi_{1,\pmb{k}}^{\dagger}) (\partial_{k_{x}} \psi_{1,\pmb{k}}) \big]
\end{equation*}
在有限温度下，我们有
\begin{equation*}
\sigma_{xy} = \frac{e^{2}}{h}\, \frac{1}{2\pi} \int \mathrm{d}^{2}\pmb{k}\; n_{F}(\epsilon_{\pmb{k}})\, \mathrm{i} \big[ (\partial_{k_{x}} \psi_{1,\pmb{k}}^{\dagger}) (\partial_{k_{y}} \psi_{1,\pmb{k}}) - (\partial_{k_{y}} \psi_{1,\pmb{k}}^{\dagger}) (\partial_{k_{x}} \psi_{1,\pmb{k}}) \big]
\end{equation*}
（注意，对于部分填充的能带，$K$ 不量子化，因为 $|\pmb{\theta}\rangle$ 不具有 $2\pi$ 周期性，见式 (4.4.9)。）

\subsection{注记：$|\pmb{\theta}\rangle$ 的周期结构与 $K$ 的量子化}

首先，让我们研究费米子基态 $|\pmb{\theta}\rangle$ 在 $\pmb{\theta}$ 空间中的周期结构。这里 $|\pmb{\theta}\rangle$ 是如下哈密顿量的基态
\begin{equation*}
H_{\pmb{\theta}} = \int \mathrm{d}^{2}\pmb{x} \left( -\frac{\hbar^{2}}{2m} \sum_{j=1,2} c^{\dagger} \Big( \partial_{j} - \mathrm{i} A_{j} - \frac{\mathrm{i} \theta_{j}}{L_{j}} \Big)^{2} c + V(c^{\dagger} c) \right)
\end{equation*}
我们注意到 $U(1)$ 规范变换
\begin{align*}
c(\pmb{x}) &\to \mathrm{e}^{2\mathrm{i}\pi x_{1} L_{1}^{-1}} c(\pmb{x})\\
(A_{1}, A_{2}) &\to \Big( A_{1} + \frac{2\pi}{L_{1}},\, A_{2} \Big)
\end{align*}
和
\begin{align*}
c(\pmb{x}) &\to \mathrm{e}^{2\mathrm{i}\pi x_{2} L_{2}^{-1}} c(\pmb{x})\\
(A_{1}, A_{2}) &\to \Big( A_{1},\, A_{2} + \frac{2\pi}{L_{2}} \Big)
\end{align*}
并不改变费米子算符 $c(\pmb{x})$ 在 $x$ 与 $y$ 方向上的周期性边界条件。然而，这些规范变换生成如下平移：$(\theta_{1}, \theta_{2}) \to (\theta_{1} + 2\pi, \theta_{2})$ 与 $(\theta_{1}, \theta_{2}) \to (\theta_{1}, \theta_{2} + 2\pi)$。或者更精确地说，我们有
\begin{align*}
H_{\theta_{1}+2\pi, \theta_{2}} &= W_{1} H_{\theta_{1}, \theta_{2}} W_{1}^{\dagger}
& W_{1} &= \mathrm{e}^{\mathrm{i} \int \mathrm{d}^{2}\pmb{x}\; 2\pi x_{1} L_{1}^{-1} c^{\dagger} c}\\
H_{\theta_{1}, \theta_{2}+2\pi} &= W_{2} H_{\theta_{1}, \theta_{2}} W_{2}^{\dagger}
& W_{2} &= \mathrm{e}^{\mathrm{i} \int \mathrm{d}^{2}\pmb{x}\; 2\pi x_{2} L_{2}^{-1} c^{\dagger} c}
\end{align*}
因此，如果 $|\pmb{\theta}\rangle$ 是 $\pmb{\theta}$ 处的多体基态，那么 $W_{1}|\pmb{\theta}\rangle$ 就是平移后的 $\pmb{\theta}$ 处的基态。于是
\begin{equation*}
\big| \theta_{1} + 2\pi, \theta_{2} \big\rangle = \mathrm{e}^{\mathrm{i} f_{1}(\pmb{\theta})} W_{1} \big| \theta_{1}, \theta_{2} \big\rangle \tag{4.4.7}
\end{equation*}
其中 $f_{1}(\pmb{\theta})$ 可任选。类似地，我们有
\begin{equation*}
\big| \theta_{1}, \theta_{2} + 2\pi \big\rangle = \mathrm{e}^{\mathrm{i} f_{2}(\pmb{\theta})} W_{2} \big| \theta_{1}, \theta_{2} \big\rangle \tag{4.4.8}
\end{equation*}

这两个幺正算符 $\mathrm{e}^{\mathrm{i} f_{1}(\pmb{\theta})} W_{1}$ 与 $\mathrm{e}^{\mathrm{i} f_{2}(\pmb{\theta})} W_{2}$ 生成把 $|\theta_{1}, \theta_{2}\rangle$、$|\theta_{1} + 2\pi, \theta_{2}\rangle$ 与 $|\theta_{1}, \theta_{2} + 2\pi\rangle$ 联系起来的那些规范变换。

我们总可以重新定义 $|\pmb{\theta}\rangle$ 的相位：$|\pmb{\theta}\rangle \to \mathrm{e}^{\mathrm{i} \phi(\pmb{\theta})} |\pmb{\theta}\rangle$，使得 $f_{1} = 0$，只要我们选择 $\phi$ 满足
\begin{equation*}
\phi(\theta_{1} + 2\pi, \theta_{2}) - \phi(\theta_{1}, \theta_{2}) = f_{1}(\theta_{1}, \theta_{2}).
\end{equation*}
那么，由
\begin{equation*}
\big| \theta_{1} + 2\pi, \theta_{2} + 2\pi \big\rangle = \mathrm{e}^{\mathrm{i} f_{2}(\theta_{1}+2\pi, \theta_{2})} W_{2} \big| \theta_{1} + 2\pi, \theta_{2} \big\rangle
\end{equation*}
我们得到
\begin{equation*}
W_{1} \big| \theta_{1}, \theta_{2} + 2\pi \big\rangle = \mathrm{e}^{\mathrm{i} f_{2}(\theta_{1}+2\pi, \theta_{2})} W_{2} W_{1} \big| \theta_{1}, \theta_{2} \big\rangle
\end{equation*}
由于 $[W_{1}, W_{2}] = 0$，我们有
\begin{equation*}
\big| \theta_{1}, \theta_{2} + 2\pi \big\rangle = \mathrm{e}^{\mathrm{i} f_{2}(\theta_{1}+2\pi, \theta_{2})} W_{2} \big| \theta_{1}, \theta_{2} \big\rangle
\end{equation*}
把上式与式 (4.4.8) 比较，我们发现
\begin{equation*}
\mathrm{e}^{\mathrm{i} f_{2}(\theta_{1}+2\pi, \theta_{2})} = \mathrm{e}^{\mathrm{i} f_{2}(\theta_{1}, \theta_{2})}
\end{equation*}
我们看到 $|\pmb{\theta}\rangle$ 在 $\theta_{1}$ 与 $\theta_{2}$ 上以 $2\pi$ 为周期准周期：
\begin{align*}
\big| \theta_{1} + 2\pi, \theta_{2} \big\rangle &= W_{1} \big| \theta_{1}, \theta_{2} \big\rangle\\
\big| \theta_{1}, \theta_{2} + 2\pi \big\rangle &= \mathrm{e}^{\mathrm{i} f_{2}(\theta_{1}, \theta_{2})} W_{2} \big| \theta_{1}, \theta_{2} \big\rangle \tag{4.4.9}
\end{align*}
其中 $W_{1,2}$ 是不依赖 $\pmb{\theta}$ 的幺正算符。

利用 $\langle \pmb{\theta} | \partial_{\pmb{\theta}} | \pmb{\theta} \rangle = \langle \pmb{\theta} | W_{1,2}^{\dagger} \partial_{\pmb{\theta}} W_{1,2} | \pmb{\theta} \rangle$ 与式 (4.4.9)，我们可以把贝里相 (4.4.3) 约化为
\begin{equation*}
2\pi K = \int \mathrm{d}\theta_{1}\; \partial_{\theta_{1}} f_{2}(\theta_{1}, 0) = f_{2}(2\pi, 0) - f_{2}(0, 0)
\end{equation*}
由于 $\mathrm{e}^{\mathrm{i} f_{2}(\theta_{1}, \theta_{2})}$ 在 $\theta_{1}$ 上是周期的，我们发现 $f_{2}(2\pi, 0) - f_{2}(0, 0) = 2\pi \times$ 整数，即 $K$ 量子化为一个整数。
